% META: {"id": "TSPE-LOG-CR-013", "chapitre": "TSPE-LOGARITHME", "type_objet": "cours", "capacites_codes": ["C4"], "sources_inspiration": [], "mode_creation": "ex_nihilo", "status": "generated"}

\section{Limite de $x \ln x$ en $0$}

% ============================================================
% STRATE 1 — Demonstration exigible
% ============================================================

\theoreme{
\[
  \lim_{x \to 0^+} x \ln x = 0.
\]
}

\demonstration{
Posons $X = -\ln x$. Quand $x \to 0^+$, $\ln x \to -\infty$, donc $X = -\ln x \to +\infty$.

De $X = -\ln x$, on tire $\ln x = -X$, soit $x = \mathrm{e}^{-X}$.

Alors :
\[
  x \ln x = \mathrm{e}^{-X} \times (-X) = -\frac{X}{\mathrm{e}^{X}}.
\]

D'apres la croissance comparee $\dfrac{\mathrm{e}^X}{X} \to +\infty$ quand $X \to +\infty$ (chapitre limites de fonctions), on a $\dfrac{X}{\mathrm{e}^X} \to 0$, donc $x \ln x = -\dfrac{X}{\mathrm{e}^X} \to 0$.
}

\margeAppui{
Ce resultat prolonge par continuite : on pose par convention $0 \ln 0 = 0$, ce qui permet de definir $x \mapsto x\ln x$ en $0$ (utile par exemple en theorie de l'information, entropie).
}

\exemple{
Calculer $\displaystyle\lim_{x \to 0^+} x^2 \ln x$.

On ecrit $x^2 \ln x = x \times (x \ln x)$. Or $x \to 0$ et $x \ln x \to 0$, donc par produit des limites, $x^2 \ln x \to 0$.

% BEGIN-VERIFY
% from sympy import symbols, ln, limit
% x = symbols('x', positive=True)
% assert limit(x*ln(x), x, 0, '+') == 0
% assert limit(x**2*ln(x), x, 0, '+') == 0
% END-VERIFY
}

% ============================================================
% STRATE 2 — Erreurs frequentes
% ============================================================

\erreurFrequente{
\textbf{Conclure trop vite sur une forme « $0 \times (-\infty)$ ».} $x \ln x$ en $0^+$ est une forme indeterminee (produit d'un terme qui tend vers $0$ et d'un terme qui tend vers $-\infty$) : il faut un changement de variable et la croissance comparee pour conclure, on ne peut pas se contenter d'un raisonnement intuitif.
}
